% =====================================================================
\chapter{El algoritmo BPE}
\label{tema:2}
% =====================================================================

El algoritmo BPE (\emph{Byte Pair Encoding}) construye un vocabulario de subpalabras de abajo arriba: parte de los símbolos básicos y fusiona, una y otra vez, el par adyacente más frecuente. Lo propuso Gage como método de compresión de datos~\cite{gage}, y Sennrich, Haddow y Birch lo adaptaron a la segmentación de palabras para la traducción automática~\cite[§3.2]{sennrich}; Radford \emph{et al.} lo aplicaron a bytes en GPT-2~\cite[§2.2]{gpt2}. Este tema define el algoritmo, demuestra cuánto acorta cada fusión y prueba que la decodificación recupera el texto. Se sigue~\cite[§2.4]{jm3} y la §2.4 de la tarea~1 de CS336~\cite{cs336}.

\section{Entrenamiento}

\begin{definicion}[Corpus de palabras y pares]\label{def:corpus-bpe}
Un \emph{corpus de entrenamiento} es un conjunto finito de palabras $w$, cada una con una frecuencia $n_w\ge 1$ y escrita como una secuencia de símbolos de un vocabulario $\V$. Para dos símbolos $a,b\in\V$, la \emph{frecuencia del par} $(a,b)$ es
\[
N(a,b)=\sum_{w}n_w\cdot\#\{\,i : w_i=a,\ w_{i+1}=b\,\},
\]
el número de posiciones en que $b$ sigue inmediatamente a $a$, ponderado por la frecuencia de cada palabra.
\end{definicion}

\begin{definicion}[Fusión]
La \emph{fusión} del par $(a,b)$ sustituye en cada palabra, recorriéndola de izquierda a derecha, cada aparición de $a$ seguido de $b$ por un símbolo nuevo $ab=a\Vert b$, y añade $ab$ al vocabulario.
\end{definicion}

\begin{definicion}[Entrenamiento de BPE]\label{def:bpe}
Dados un corpus escrito en un vocabulario inicial $\V_0$ y un número de fusiones $k$, el \emph{entrenamiento de BPE} repite $k$ veces:
\begin{enumerate}[label=\arabic*.]
  \item calcular $N(a,b)$ para todos los pares adyacentes;
  \item elegir un par $(a^{*},b^{*})$ de frecuencia máxima, con una regla de desempate fija;
  \item fusionarlo y anotar la regla $a^{*}+b^{*}\to a^{*}b^{*}$.
\end{enumerate}
El resultado es la \emph{lista ordenada de fusiones} $\mu_1,\dots,\mu_k$, que constituye el tokenizador.
\end{definicion}

Las palabras se obtienen del texto en un paso previo, la \emph{pretokenización}, de modo que las fusiones nunca cruzan el límite entre dos palabras~\cite[§3.2]{sennrich}. Sennrich \emph{et al.} añaden al final de cada palabra un símbolo especial de fin de palabra para poder reconstruir los espacios.

\begin{proposicion}\label{prop:tam-vocab}
Tras $k$ fusiones, el vocabulario tiene a lo sumo $\abs{\V_0}+k$ símbolos, con igualdad cuando las $k$ cadenas creadas son distintas entre sí y no pertenecen a $\V_0$. En particular, si $\V_0$ está formado por símbolos de longitud~1, $\abs{\V}=\abs{\V_0}+k$ salvo repeticiones de una misma cadena obtenida por dos caminos distintos.
\end{proposicion}

En la práctica se usa la igualdad $\abs{\V}=\abs{\V_0}+k$, y $k$ es el único hiperparámetro del entrenamiento~\cite[§3.2]{sennrich}.

\begin{lema}\label{lem:solape}
Si $a\neq b$, dos apariciones del par $(a,b)$ en una secuencia no comparten ninguna posición.
\end{lema}

\begin{proposicion}[Compresión de una fusión]\label{prop:compresion}
Sea un corpus con $L$ símbolos en total (contando cada palabra $n_w$ veces) y sea $(a,b)$, con $a\neq b$, un par de frecuencia $N(a,b)=c$. Tras fusionarlo, el corpus tiene $L-c$ símbolos.
\end{proposicion}

\begin{observacion}
Si $a=b$, las apariciones pueden solaparse: «\tok{aaa}» contiene dos pares $(\tok{a},\tok{a})$, pero la fusión de izquierda a derecha solo puede aplicar una y deja «\tok{aa}\sep\tok{a}». En ese caso la longitud baja en el número de apariciones \emph{no solapadas}. Por la \cref{prop:compresion}, elegir el par más frecuente es elegir la fusión que más acorta el corpus en ese paso: BPE es un algoritmo \emph{voraz}. No garantiza la segmentación globalmente más corta con $k$ fusiones, pero funciona bien en la práctica.
\end{observacion}

\addtocounter{definicion}{1}%

\begin{table}[h]
\centering\footnotesize
\setlength{\tabcolsep}{4pt}
\caption{Cuatro fusiones de BPE sobre el corpus del \cref{ej:bpe-casa}. El punto separa los símbolos de cada palabra.}\label{tab:bpe-casa}
\begin{tabularx}{\textwidth}{@{}cll>{\raggedright\arraybackslash}Xr@{}}
\toprule
Paso & Pares más frecuentes & Fusión & Corpus resultante & $L$ \\
\midrule
0 & --- & --- & \tok{c}\sep\tok{a}\sep\tok{s}\sep\tok{a}\,($\times3$); \tok{c}\sep\tok{a}\sep\tok{s}\sep\tok{a}\sep\tok{s}\,($\times2$); \tok{m}\sep\tok{a}\sep\tok{s}\sep\tok{a}; \tok{m}\sep\tok{e}\sep\tok{s}\sep\tok{a} & 30 \\
1 & $(\tok{a},\tok{s})$ 8, $(\tok{s},\tok{a})$ 7 & $\tok{a}+\tok{s}$ & \tok{c}\sep\tok{as}\sep\tok{a}\,($\times3$); \tok{c}\sep\tok{as}\sep\tok{as}\,($\times2$); \tok{m}\sep\tok{as}\sep\tok{a}; \tok{m}\sep\tok{e}\sep\tok{s}\sep\tok{a} & 22 \\
2 & $(\tok{c},\tok{as})$ 5, $(\tok{as},\tok{a})$ 4 & $\tok{c}+\tok{as}$ & \tok{cas}\sep\tok{a}\,($\times3$); \tok{cas}\sep\tok{as}\,($\times2$); \tok{m}\sep\tok{as}\sep\tok{a}; \tok{m}\sep\tok{e}\sep\tok{s}\sep\tok{a} & 17 \\
3 & $(\tok{cas},\tok{a})$ 3, $(\tok{cas},\tok{as})$ 2 & $\tok{cas}+\tok{a}$ & \tok{casa}\,($\times3$); \tok{cas}\sep\tok{as}\,($\times2$); \tok{m}\sep\tok{as}\sep\tok{a}; \tok{m}\sep\tok{e}\sep\tok{s}\sep\tok{a} & 14 \\
4 & $(\tok{cas},\tok{as})$ 2, resto 1 & $\tok{cas}+\tok{as}$ & \tok{casa}\,($\times3$); \tok{casas}\,($\times2$); \tok{m}\sep\tok{as}\sep\tok{a}; \tok{m}\sep\tok{e}\sep\tok{s}\sep\tok{a} & 12 \\
\bottomrule
\end{tabularx}
\end{table}

\section{Codificación y decodificación}

\begin{definicion}[Codificación BPE]\label{def:cod-bpe}
Para codificar una palabra nueva se escribe en símbolos de $\V_0$ y se aplican las fusiones $\mu_1,\dots,\mu_k$ \emph{en el orden en que se aprendieron}, cada una a toda la palabra, de izquierda a derecha. La \emph{decodificación} de una secuencia de tokens es la concatenación de las cadenas que representan.
\end{definicion}

\addtocounter{definicion}{1}%

\begin{teorema}[Corrección de BPE]\label{teo:bpe-correcto}
Sea $\V_0$ un vocabulario inicial tal que toda cadena admite una escritura en símbolos de $\V_0$, y sea $\delta$ la decodificación por concatenación. Para toda cadena $x$, la codificación $\tau(x)$ está definida y $\delta(\tau(x))=x$.
\end{teorema}

\begin{corolario}\label{cor:no-unk}
Un tokenizador BPE nunca falla: en el peor caso no se aplica ninguna fusión y se devuelven los símbolos de $\V_0$. Con $\V_0=\mathcal{B}$ (BPE a nivel de byte), por el \cref{cor:vocab256} no hay tokens desconocidos.
\end{corolario}

\section{Desempates y determinismo}

Cuando varios pares empatan en frecuencia máxima, el paso~2 de la \cref{def:bpe} necesita una regla. Cualquier regla sirve, pero debe ser fija para que el tokenizador sea reproducible. La tarea de CS336 elige el par mayor en orden lexicográfico~\cite[§2.4]{cs336}; otra opción frecuente es el primer par que aparece en el corpus. Reglas distintas producen tokenizadores distintos (\cref{eje:2-lex}).

\begin{observacion}
La codificación de la \cref{def:cod-bpe} es determinista: la misma cadena da siempre la misma secuencia de tokens. El modelo unigram del \cref{tema:3}, en cambio, define una distribución sobre segmentaciones y permite muestrearlas.
\end{observacion}

\section{BPE a nivel de byte y pretokenización}

GPT-2 aplica BPE a los bytes UTF-8 del texto, con $\V_0=\mathcal{B}$ y $\abs{\V_0}=256$, en lugar de los más de $1{,}3\cdot 10^{5}$ símbolos que exigiría partir de los puntos de código~\cite[§2.2]{gpt2}. Su vocabulario tiene $50\,257$ tokens~\cite[§2.3]{gpt2}: los 256 bytes, un token especial de fin de texto y $50\,000$ fusiones (\cref{eje:2-fusiones}).

\begin{definicion}[Pretokenización]
La \emph{pretokenización} divide el texto en fragmentos mediante una expresión regular antes de aplicar BPE, y las fusiones se aplican dentro de cada fragmento~\cite[§2.6.9]{jm3}. En GPT-2, cada fragmento es típicamente una palabra con su espacio delante (\tok{\vs lugar}, \tok{\vs de}), un signo de puntuación o un grupo de dígitos; así se impide fusionar bytes de categorías distintas~\cite[§2.2]{gpt2}.
\end{definicion}

Sin pretokenización, BPE aprende tokens que mezclan palabras y puntuación, como «\tok{de\vs su\vs hacienda.\vs}». Una consecuencia de la pretokenización es que \tok{\vs hidalgo} y \tok{hidalgo} son tokens distintos: un espacio de más al final de un texto cambia su tokenización.

\begin{observacion}[Tokens que contienen medio carácter]
Con $\V_0=\mathcal{B}$, un token puede ser un fragmento de un carácter multibyte, por ejemplo el byte \hx{C3} solo. Por el \cref{teo:bpe-correcto}, la concatenación de todos los tokens de un texto es su cadena UTF-8, que se decodifica sin pérdida (\cref{teo:prefijo}). Pero un token aislado puede no ser UTF-8 válido, y una interfaz que decodifica cada token al recibirlo debe guardar en un búfer los bytes de un carácter incompleto (\cref{eje:2-medio}).
\end{observacion}

